% ============================================================
\chapter{Variables Al\'eatoires Continues}
\label{ch:continues}
% ============================================================

Avec les variables discr\`etes, on comptait. Avec les variables continues, on mesure. La transition est plus qu'un changement de vocabulaire~: elle exige de remplacer les sommes par des int\'egrales, les masses ponctuelles par des densit\'es, et d'apprivoiser l'id\'ee qu'en continu, la probabilit\'e d'obtenir une valeur \emph{exacte} est toujours nulle. C'est Abraham de Moivre qui, vers 1733, a ouvert cette voie en approchant la loi binomiale par la courbe en cloche --- ce que nous appelons aujourd'hui la loi normale, et que Gauss immortalisera dans ses travaux d'astronomie.

\begin{intuition}
Une variable al\'eatoire continue prend ses valeurs dans un intervalle (ou tout
$\R$) et sa loi est d\'ecrite par une \textbf{densit\'e de probabilit\'e}. Ce
chapitre d\'eveloppe les outils sp\'ecifiques au cas continu~: densit\'e, fonction
de r\'epartition, quantiles.
\end{intuition}

% ------------------------------------------------------------
\section{Densit\'e de probabilit\'e}
% ------------------------------------------------------------

\begin{definition}[Variable al\'eatoire \`a densit\'e]
Une variable al\'eatoire $X$ est dite \textbf{continue} (ou \`a densit\'e) s'il
existe une fonction $f_X : \R \to [0, +\infty)$, int\'egrable, telle que pour tout
$a \leq b$~:
\[
    \PP(a \leq X \leq b) = \int_a^b f_X(x)\, dx.
\]
La fonction $f_X$ est appel\'ee \textbf{densit\'e de probabilit\'e} de $X$.
\end{definition}

\begin{proposition}[Propri\'et\'es de la densit\'e]
\begin{enumerate}[label=(\roman*)]
    \item $f_X(x) \geq 0$ pour tout $x \in \R$.
    \item $\int_{-\infty}^{+\infty} f_X(x)\, dx = 1$.
    \item $\PP(X = a) = 0$ pour tout $a \in \R$.
    \item $\PP(a \leq X \leq b) = \PP(a < X < b) = \PP(a \leq X < b) = \PP(a < X \leq b)$.
\end{enumerate}
\end{proposition}

\begin{attention}[title=\textbf{La densit\'e n'est pas une probabilit\'e}]
On peut avoir $f_X(x) > 1$ pour certaines valeurs de $x$. La densit\'e est une
\og concentration de probabilit\'e \fg, pas une probabilit\'e. Seule son
int\'egrale sur un intervalle donne une probabilit\'e.
\end{attention}

% ------------------------------------------------------------
\section{Fonction de r\'epartition}
% ------------------------------------------------------------

\begin{definition}[Fonction de r\'epartition]
\[
    F_X(x) = \PP(X \leq x) = \int_{-\infty}^{x} f_X(t)\, dt, \quad x \in \R.
\]
\end{definition}

\begin{theoreme}[Lien densit\'e--r\'epartition]
Si $F_X$ est d\'erivable en $x$, alors $f_X(x) = F_X'(x)$. Plus g\'en\'eralement,
$f_X$ est une d\'eriv\'ee de $F_X$ presque partout.
\end{theoreme}

\begin{proposition}[Propri\'et\'es de $F_X$]
\begin{enumerate}[label=(\roman*)]
    \item $F_X$ est croissante et continue (dans le cas \`a densit\'e).
    \item $\lim_{x \to -\infty} F_X(x) = 0$, $\lim_{x \to +\infty} F_X(x) = 1$.
    \item $\PP(a < X \leq b) = F_X(b) - F_X(a)$.
\end{enumerate}
\end{proposition}

% ------------------------------------------------------------
\section{Lois continues classiques}
% ------------------------------------------------------------

\subsection{Loi uniforme continue}

\begin{definition}[Uniforme sur $\lbrack a,b\rbrack$]
$X \sim \mathcal{U}([a,b])$ avec $a < b$~:
\[
    f_X(x) = \frac{1}{b-a}\, \mathbf{1}_{[a,b]}(x), \quad
    F_X(x) = \begin{cases} 0 & \text{si } x < a, \\ \frac{x-a}{b-a} & \text{si } a \leq x \leq b, \\ 1 & \text{si } x > b. \end{cases}
\]
$\E[X] = (a+b)/2$, $\Var(X) = (b-a)^2/12$.
\end{definition}

\subsection{Loi exponentielle}

\begin{definition}[Exponentielle]
$X \sim \text{Exp}(\lambda)$ avec $\lambda > 0$~:
\[
    f_X(x) = \lambda e^{-\lambda x}\, \mathbf{1}_{[0,+\infty)}(x), \quad
    F_X(x) = (1 - e^{-\lambda x})\, \mathbf{1}_{[0,+\infty)}(x).
\]
$\E[X] = 1/\lambda$, $\Var(X) = 1/\lambda^2$.
\end{definition}

\begin{theoreme}[Perte de m\'emoire]
La loi exponentielle est la seule loi continue sur $[0,+\infty)$ v\'erifiant~:
\[
    \PP(X > s + t \mid X > s) = \PP(X > t) \quad \forall s, t \geq 0.
\]
\end{theoreme}

\begin{proof}
Soit $g(t) = \PP(X > t) = e^{-\lambda t}$. Alors~:
\[
    \PP(X > s+t \mid X > s) = \frac{e^{-\lambda(s+t)}}{e^{-\lambda s}} = e^{-\lambda t} = g(t).
\]
R\'eciproquement, si $g(s+t) = g(s)g(t)$ avec $g$ continue d\'ecroissante et $g(0)=1$,
alors $g(t) = e^{-\lambda t}$ pour un certain $\lambda > 0$, ce qui caract\'erise
la loi exponentielle.
\end{proof}

\begin{intuition}[title=\textbf{Exponentielle et Poisson}]
Si les arriv\'ees suivent un processus de Poisson de taux $\lambda$, alors le temps
entre deux arriv\'ees cons\'ecutives suit une loi $\text{Exp}(\lambda)$.
\end{intuition}

\subsection{Loi normale (gaussienne)}

\begin{definition}[Loi normale]
$X \sim \mathcal{N}(\mu, \sigma^2)$ avec $\mu \in \R$ et $\sigma > 0$~:
\[
    f_X(x) = \frac{1}{\sigma\sqrt{2\pi}}
    \exp\!\left(-\frac{(x-\mu)^2}{2\sigma^2}\right), \quad x \in \R.
\]
$\E[X] = \mu$, $\Var(X) = \sigma^2$.
\end{definition}

\begin{definition}[Loi normale centr\'ee r\'eduite]
Si $Z \sim \mathcal{N}(0,1)$, on note $\Phi(x) = \PP(Z \leq x)$ sa fonction de
r\'epartition. On a la relation~:
\[
    \text{si } X \sim \mathcal{N}(\mu, \sigma^2), \quad
    \PP(X \leq x) = \Phi\!\left(\frac{x - \mu}{\sigma}\right).
\]
\end{definition}

\begin{formules}[title=\textbf{Valeurs cl\'es de $\Phi$}]
\begin{center}
\renewcommand{\arraystretch}{1.3}
\begin{tabular}{cccccc}
    \toprule
    $z$ & 1{,}645 & 1{,}96 & 2{,}326 & 2{,}576 & 3{,}291 \\
    \midrule
    $\Phi(z)$ & 0{,}95 & 0{,}975 & 0{,}99 & 0{,}995 & 0{,}9995 \\
    \bottomrule
\end{tabular}
\end{center}
Par sym\'etrie~: $\Phi(-z) = 1 - \Phi(z)$.
\end{formules}

\begin{proposition}[R\`egle des $k\sigma$]
Si $X \sim \mathcal{N}(\mu, \sigma^2)$~:
\begin{itemize}
    \item $\PP(\abs{X - \mu} \leq \sigma) \approx 0{,}6827$ \quad (r\`egle du $1\sigma$),
    \item $\PP(\abs{X - \mu} \leq 2\sigma) \approx 0{,}9545$ \quad (r\`egle du $2\sigma$),
    \item $\PP(\abs{X - \mu} \leq 3\sigma) \approx 0{,}9973$ \quad (r\`egle du $3\sigma$).
\end{itemize}
\end{proposition}

\subsection{Loi gamma}

\begin{definition}[Gamma]
$X \sim \Gamma(\alpha, \lambda)$ avec $\alpha, \lambda > 0$~:
\[
    f_X(x) = \frac{\lambda^\alpha}{\Gamma(\alpha)}\, x^{\alpha - 1}\, e^{-\lambda x}\,
    \mathbf{1}_{(0,+\infty)}(x),
\]
o\`u $\Gamma(\alpha) = \int_0^{\infty} t^{\alpha-1} e^{-t}\, dt$. On a
$\E[X] = \alpha/\lambda$ et $\Var(X) = \alpha/\lambda^2$.
\end{definition}

\begin{remarque}
$\Gamma(1, \lambda) = \text{Exp}(\lambda)$ et $\Gamma(n/2, 1/2) = \chi^2(n)$
(loi du chi-deux \`a $n$ degr\'es de libert\'e).
\end{remarque}

\subsection{Loi b\^eta}

\begin{definition}[B\^eta]
$X \sim \text{Beta}(\alpha, \beta)$ avec $\alpha, \beta > 0$~:
\[
    f_X(x) = \frac{x^{\alpha-1}(1-x)^{\beta-1}}{B(\alpha,\beta)}\,
    \mathbf{1}_{(0,1)}(x),
\]
o\`u $B(\alpha,\beta) = \Gamma(\alpha)\Gamma(\beta)/\Gamma(\alpha+\beta)$.
$\E[X] = \alpha/(\alpha+\beta)$.
\end{definition}

% ------------------------------------------------------------
\section{Quantiles et m\'ediane}
% ------------------------------------------------------------

\begin{definition}[Quantile]
Le \textbf{quantile d'ordre $p$} ($0 < p < 1$) de $X$ est~:
\[
    q_p = \inf\{x \in \R : F_X(x) \geq p\} = F_X^{-1}(p).
\]
\end{definition}

\begin{definition}[M\'ediane]
La \textbf{m\'ediane} de $X$ est le quantile d'ordre $1/2$~:
$\text{Med}(X) = q_{1/2}$.
\end{definition}

\begin{exemple}
Si $X \sim \text{Exp}(\lambda)$, alors $F_X(x) = 1 - e^{-\lambda x}$ pour $x \geq 0$.
Le quantile d'ordre $p$ est~:
\[
    q_p = -\frac{\ln(1-p)}{\lambda}.
\]
En particulier, $\text{Med}(X) = \ln 2 / \lambda$.
\end{exemple}

% ------------------------------------------------------------
\section{Simulation par inversion}
% ------------------------------------------------------------

\begin{theoreme}[M\'ethode d'inversion]
Si $U \sim \mathcal{U}([0,1])$ et $F$ est une fonction de r\'epartition continue
strictement croissante, alors $X = F^{-1}(U)$ a pour fonction de r\'epartition $F$.
\end{theoreme}

\begin{proof}
$\PP(X \leq x) = \PP(F^{-1}(U) \leq x) = \PP(U \leq F(x)) = F(x)$.
\end{proof}

\begin{exemple}
Pour simuler $X \sim \text{Exp}(\lambda)$, on g\'en\`ere $U \sim \mathcal{U}([0,1])$
et on pose $X = -\ln(1-U)/\lambda = -\ln(U)/\lambda$ (car $1-U$ a la m\^eme loi que $U$).
\end{exemple}

% ------------------------------------------------------------
\section{Exercices}
% ------------------------------------------------------------

\begin{exercice}
Soit $X$ de densit\'e $f(x) = cx^2\, \mathbf{1}_{[0,1]}(x)$. D\'eterminer $c$,
$\E[X]$, $\Var(X)$ et $F_X$.
\end{exercice}

\begin{exercice}
Soit $X \sim \mathcal{N}(100, 25)$. Calculer~:
\begin{enumerate}[label=(\alph*)]
    \item $\PP(X > 105)$\,;
    \item $\PP(90 < X < 110)$\,;
    \item le quantile $q_{0{,}9}$.
\end{enumerate}
\end{exercice}

\begin{exercice}
Montrer que si $X \sim \mathcal{N}(\mu, \sigma^2)$, alors $Z = (X - \mu)/\sigma
\sim \mathcal{N}(0,1)$.
\end{exercice}

\begin{exercice}
Soit $X \sim \text{Exp}(\lambda)$. Calculer la densit\'e et l'esp\'erance de
$Y = X^2$.
\end{exercice}

\begin{exercice}
Montrer que l'int\'egrale de Gauss vaut~:
$\int_{-\infty}^{+\infty} e^{-x^2/2}\, dx = \sqrt{2\pi}$.
\textit{Indication~:} calculer $I^2$ en coordonn\'ees polaires.
\end{exercice}

\begin{exercice}
Soit $X \sim \Gamma(\alpha, \lambda)$. Montrer que $Y = 2\lambda X \sim \chi^2(2\alpha)$
lorsque $\alpha \in \N^*/2$.
\end{exercice}

\begin{exercice}[Simulation]
D\'ecrire une m\'ethode pour simuler une variable al\'eatoire de densit\'e
$f(x) = 2x\, \mathbf{1}_{[0,1]}(x)$ par la m\'ethode d'inversion.
\end{exercice}

\begin{exercice}
Soit $X$ de densit\'e $f(x) = \frac{2}{\pi(1+x^2)}\, \mathbf{1}_{(0,+\infty)}(x)$.
V\'erifier que $f$ est bien une densit\'e et calculer $F_X$.
L'esp\'erance de $X$ existe-t-elle~?
\end{exercice}
